%% kapa: exam=1 round=33 year=2022 period=1 form=A subject=civil_law subject_ko=민법 seq=17 number=17
%% source: Q-Net 감정평가사 기출문제 (cst00309) artlSeq=5209997 file=raw/1008/5209997/2022년도 제33회 감정평가사 1차시험 1교시(민법,경제학,부동산학).pdf
%% conversion: kordoc 4.24.51 -> tools/gen1.py (auto); review=proofread-sonnet; flags=none
\begin{question}{17}
법률행위의 취소에 관한 설명으로 옳지 않은 것은? (다툼이 있으면 판례에 따름)
\begin{choices}
  \item 취소할 수 있는 미성년자의 법률행위를 친권자가 추인하는 경우, 그 취소의 원인이 소멸한 후에 하여야만 효력이 있다.
  \item 제한능력자가 그 의사표시를 취소한 경우, 제한능력자는 그 행위로 인하여 받은 이익이 현존하는 한도에서 상환(償還)할 책임이 있다.
  \item 강박에 의하여 의사표시를 한 자의 포괄승계인은 그 의사표시를 취소할 수 있다.
  \item 취소권은 추인할 수 있는 날로부터 3년 내에, 법률행위를 한 날로부터 10년 내에 행사하여야 한다.
  \item 의사표시의 취소는 취소기간 내에 소를 제기하는 방법으로만 행사하여야 하는 것은 아니다.
\end{choices}
\end{question}
\answer{1}
